\chapter{Conclusion}
\label{cap:conclusao}

This work investigated how data leakage introduced at four steps of an ML
pipeline (normalization, missing value imputation, feature selection, and
hyperparameter tuning) distorts the performance estimates of regression models.
Following a paired design, leakage-free and leakage-prone pipelines were run on
the same data partitions for 21 regression data sets from PMLB and five learning
algorithms, totaling 11,340 paired comparisons, and the differences between them
were measured with $R^2$ and RMSE and analyzed with non-parametric tests.

The results show that the danger of leakage depends above all on \emph{what}
leaks. Fitting a scaler or an imputer on the full data set leaks only statistics
of the features, and its effect on the estimates was negligible: mean
differences in $R^2$ of about $2\times10^{-4}$, distributed almost symmetrically
around zero, with the single exception of a small but significant bias for SVR
under imputation. Steps that use the target are different. Non-nested
hyperparameter tuning inflated the estimates of all five algorithms in a
statistically significant way, with differences in $R^2$ of up to 0.35 and RMSE
underestimated by up to 54\%. Feature selection performed before data splitting
left most configurations unchanged, but produced some of the largest distortions
of the study, such as a model whose $R^2$ rose from $-0.13$ to $0.49$ when the
selection saw the test instances.

The study also shows that the effects are concentrated in small data sets. For
hyperparameter tuning, the inflation was strongly correlated with the number of
instances and with the ratio between features and instances, and in all
scenarios the largest distortions occurred in data sets with fewer than 1,000
instances. Regarding the algorithms, differences were significant only for
tuning, where KNN was the most consistently affected and LR and SVR showed the
most extreme individual cases; tree-based models, and RF in particular, were the
most robust overall.

From a practical standpoint, these findings reinforce well-known recommendations
and help prioritize them. Every step that learns from data should be part of the
pipeline that is cross-validated, and this is mandatory for any step that uses
the target, such as supervised feature selection, and for hyperparameter
tuning, whose performance must be estimated with nested CV or with a test set
that took no part in the search. The smaller the data set, the more critical
these precautions become, and the less a single optimistic number obtained from a
leaky protocol should be trusted. Conversely, the small effects observed for
normalization and imputation do not make these mistakes acceptable, but suggest
that, in data sets of moderate size and under the conditions tested here, they
are unlikely to change the conclusions of a study.

\section{Future Work}

Several extensions could deepen these findings. The imputation scenario could be
repeated with a larger fraction of incomplete features and with MCAR and MNAR
mechanisms, which would test whether the small effects observed here persist
under heavier missingness. Feature selection could be evaluated with wrapper and
embedded methods, which interact more closely with the learning algorithm, and
with a larger number of small and high-dimensional data sets, where this form of
leakage is expected to matter most. Combining several leakage sources in the same
pipeline would reveal whether their effects add up, compensate, or amplify each
other. Other algorithms, such as neural networks and gradient boosting, and other
search strategies, such as random search and Bayesian optimization, could also be
included. Finally, the relation between data set size and leakage suggests that a
controlled study varying the number of instances of the same data sets, through
subsampling, could provide a more precise characterization of when leakage
becomes harmful in regression.
